% PDF page 157
\source{gilbarg-trudinger:page-0157}


\noindent\textbf{7.1.} \textit{$L^p$ Spaces}\label{gt:sec-lp-spaces}\dcref{7.1}\group{ch07-sobolev-spaces}\source{gilbarg-trudinger:page-0157}

\medskip

Throughout this chapter $\Omega$ will denote a bounded domain in $\mathbb{R}^n$. By a measurable function on $\Omega$ we shall mean an equivalence class of measurable functions on $\Omega$ which differ only on a subset of measure zero. Any pointwise property attributed to a measurable function will thus be understood to hold in the usual sense for some function in the same equivalence class. The supremum and infimum of a measurable function will then be understood as the essential supremum and infimum.

For $p \geq 1$, we let $L^p(\Omega)$ denote the classical Banach space consisting of measurable functions on $\Omega$ that are $p$-integrable. The norm in $L^p(\Omega)$ is defined by

\begin{equation}
\tag{7.3}
\|u\|_{p;\Omega} = \|u\|_{L^p(\Omega)} = \left(\int_\Omega |u|^p\,dx\right)^{1/p}.
\end{equation}

When $u$ is a vector or matrix function the same notation will be used, the norm $|u|$ denoting the usual Euclidean norm. For $p=\infty$, $L^\infty(\Omega)$ denotes the Banach space of bounded functions on $\Omega$ with the norm

\begin{equation}
\tag{7.4}
\|u\|_{\infty;\Omega} = \|u\|_{L^\infty(\Omega)} = \sup_\Omega |u|.
\end{equation}

In the following we shall use $\|u\|_p$ for $\|u\|_{L^p(\Omega)}$ when there is no ambiguity.

We shall need the following inequalities in dealing with integral estimates:

\textit{Young's inequality.}

\begin{equation}
\tag{7.5}
ab \leq \frac{a^p}{p} + \frac{b^q}{q};
\end{equation}

this holds for positive real numbers $a, b, p, q$ satisfying

\[
\frac{1}{p} + \frac{1}{q} = 1.
\]

The case $p=q=2$ of inequality (7.5) is known as Cauchy's inequality. Replacing $a$ by $\varepsilon^{1/p}a$, $b$ by $\varepsilon^{-1/q}b$ for positive $\varepsilon$, we obtain a useful interpolation inequality

\begin{equation}
\tag{7.6}
ab \leq \frac{\varepsilon a^p}{p} + \frac{\varepsilon^{-q/p}b^q}{q}
\leq \varepsilon a^p + \varepsilon^{-q/p}b^q.
\end{equation}

\textit{H\"older's inequality.}

\begin{equation}
\tag{7.7}
\int_\Omega uv\,dx \leq \|u\|_p \|v\|_q;
\end{equation}

this holds for functions $u \in L^p(\Omega)$, $v \in L^q(\Omega)$, $1/p + 1/q = 1$ and is a consequence of Young's inequality. When $p=q=2$, H\"older's inequality reduces to the well


% PDF page 158
\source{gilbarg-trudinger:page-0158}

known Schwarz inequality. That the expression (7.3) defines a norm on $L^p(\Omega)$
is a consequence of H\"older's inequality. Let us note some other simple consequences
of H\"older's inequality.

\begin{equation}
|\Omega|^{-1/p}\|u\|_p \le |\Omega|^{-1/q}\|u\|_q \qquad \text{for $u \in L^q(\Omega)$, \ $p \le q$.}
\tag{7.8}
\end{equation}

\begin{equation}
\|u\|_q \le \|u\|_p^\lambda \|u\|_r^{1-\lambda} \qquad \text{for $u \in L^r(\Omega)$,}
\tag{7.9}
\end{equation}

where $p \le q \le r$ and $1/q = \lambda/p + (1-\lambda)/r$.
Combining inequalities (7.6) and (7.9), we obtain an interpolation inequality for $L^p$
norms, namely,

\begin{equation}
\|u\|_q \le \varepsilon\|u\|_r + \varepsilon^{-\mu}\|u\|_p,
\tag{7.10}
\end{equation}

where

\[
\mu = \left(\frac{1}{p} - \frac{1}{q}\right)\bigg/\left(\frac{1}{q} - \frac{1}{r}\right).
\]

We shall also have occasion to use a generalization of H\"older's inequality to $m$
functions, $u_1,\ldots,u_m$, lying respectively in spaces $L^{p_1},\ldots,L^{p_m}$, where

\[
\frac{1}{p_1} + \cdots + \frac{1}{p_m} = 1.
\]

The resulting inequality, obtainable from the case $m = 2$ by an induction argument,
is then

\begin{equation}
\int_{\Omega} u_1 \cdots u_m \, dx \le \|u_1\|_{p_1}\cdots\|u_m\|_{p_m}.
\tag{7.11}
\end{equation}

It is also of interest to study the $L^p$ norm as a function of $p$. Writing

\begin{equation}
\Phi_p(u) = \left(\frac{1}{|\Omega|}\int_{\Omega} |u|^p \, dx\right)^{1/p}
\tag{7.12}
\end{equation}

for $p > 0$, we see that $\Phi$ is non-decreasing in $p$ for fixed $u$, by inequality (7.8),
while the inequality (7.9) shows that $\Phi$ is logarithmically convex in $p^{-1}$. Note that
$\Phi_p(u) = |\Omega|^{-1/p}\|u\|_p$ for $p \ge 1$. Although the functional $\Phi$ does not extend the $L^p$
norm as a norm for values of $p$ less than one, it will nevertheless be useful for later
purposes (see Chapter 8).
We also note here some of the well known functional analytic properties of the
$L^p$ spaces ; (see for example Royden [RY]). The space $L^p(\Omega)$ is separable for
$p < \infty$, $C^0(\overline{\Omega})$ being in particular a dense subspace. The dual space of $L^p(\Omega)$ is
isomorphic to $L^q(\Omega)$ provided $1/p + 1/q = 1$ and $p < \infty$. Hence $L^p(\Omega)$ is reflexive

% PDF page 159
\source{gilbarg-trudinger:page-0159}

for $1 < p < \infty$. The number $q$, the Hölder conjugate of $p$, will often be denoted $p'$. Finally, $L^2(\Omega)$ is a Hilbert space under the scalar product
\[
(u,v) = \int_{\Omega} uv\,dx.
\]

\section*{7.2. \ Regularization and Approximation by Smooth Functions}

The spaces $C^{k,\alpha}(\Omega)$ which were introduced in Chapter 4 are local spaces. Let us define local analogues of the $L^p(\Omega)$ spaces by letting $L^p_{\mathrm{loc}}(\Omega)$ denote the linear space of measurable functions locally $p$-integrable in $\Omega$. Although they are not normed spaces, the $L^p_{\mathrm{loc}}(\Omega)$ spaces are readily topologized. Namely, a sequence $\{u_m\}$ converges to $u$ in the sense of $L^p_{\mathrm{loc}}(\Omega)$ if $\{u_m\}$ converges to $u$ in $L^p(\Omega')$ for each $\Omega' \subset\subset \Omega$.

Let $\rho$ be a non-negative function in $C^\infty(\R^n)$ vanishing outside the unit ball $B_1(0)$ and satisfying
\[
\int \rho\,dx = 1.
\]
Such a function is often called a mollifier. A typical example is the function $\rho$ given by
\[
\rho(x)=
\begin{cases}
c \exp\!\left(\dfrac{1}{|x|^2-1}\right) & \text{for } |x|\le 1 \\
0 & \text{for } |x|\ge 1
\end{cases}
\]
where $c$ is chosen so that $\int \rho\,dx = 1$ and whose graph has the familiar bell shape.

For $u\in L^1_{\mathrm{loc}}(\Omega)$ and $h>0$, the regularization of $u$, denoted by $u_h$, is then defined by the convolution

\[
\tag{7.13}
u_h(x)=h^{-n}\int_{\Omega}\rho\!\left(\frac{x-y}{h}\right)u(y)\,dy
\]

provided $h<\dist(x,\partial\Omega)$. It is clear that $u_h$ belongs to $C^\infty(\Omega')$ for any $\Omega' \subset\subset \Omega$ provided $h<\dist(\Omega',\partial\Omega)$. Furthermore, if $u$ belongs to $L^1(\Omega)$, $\Omega$ bounded, then $u_h$ lies in $C^\infty(\R^n)$ for arbitrary $h>0$. As $h$ tends to zero, the function $y\mapsto h^{-n}\rho(x-y/h)$ tends to the Dirac delta distribution at the point $x$. The significant feature of regularization, which we partly explore now, is the sense in which $u_h$ approximates $u$ as $h\to 0$. It turns out, roughly stated, that if $u$ lies in a local space, then $u_h$ approximates u in the natural topology of that space.

\begin{lemma}\label{gt:lem-sobolev-regularization-uniform}\dcref{Lemma 7.1}\group{ch07-sobolev-spaces}\source{gilbarg-trudinger:page-0159}
\emph{Let $u\in C^0(\Omega)$. Then $u_h$ converges to $u$ uniformly on any domain $\Omega' \subset\subset \Omega$.}
\end{lemma}

% PDF page 160
\source{gilbarg-trudinger:page-0160}


\textit{Proof.} We have

\[
u_h(x)=h^{-n}\int_{|x-y|\le h}\rho\!\left(\frac{x-y}{h}\right)u(y)\,dy
=\int_{|z|\le 1}\rho(z)u(x-hz)\,dz
\qquad \left(\text{putting } z=\frac{x-y}{h}\right);
\]

hence if $\Omega' \Subset \Omega$ and $2h<\dist(\Omega',\partial\Omega)$,

\[
\sup_{\Omega'}|u-u_h|
\le \sup_{x\in\Omega'}\int_{|z|\le 1}\rho(z)\lvert u(x)-u(x-hz)\rvert\,dz
\]

\[
\le \sup_{x\in\Omega'}\sup_{|z|\le 1}\lvert u(x)-u(x-hz)\rvert.
\]

Since $u$ is uniformly continuous over the set

\[
B_h(\Omega')=\{x \mid \dist(x,\Omega')<h\},
\]

$u_h$ tends to $u$ uniformly on $\Omega'$. $\square$

The convergence in Lemma 7.1 would be uniform over all of $\Omega$ if $u$ vanished continuously on $\partial\Omega$. More generally if $u\in C^0(\overline{\Omega})$ we can define an extension $\tilde u$ of $u$ such that $\tilde u=u$ in $\Omega$ and $\tilde u\in C^0(\tilde\Omega)$ for some $\tilde\Omega \supset\supset \Omega$. Then $\tilde u_h$, the regularization of $\tilde u$ in $\tilde\Omega$, converges to $u$ uniformly in $\Omega$ as $h\to 0$.

The process of regularization can also be used to approximate Hölder continuous functions. In particular if $u\in C^{\alpha}(\Omega)$, $0\le \alpha\le 1$, then

\begin{equation}
\tag{7.14}
[u_h]_{\alpha;\Omega'}\le [u]_{\alpha;\Omega''}
\end{equation}

where $\Omega''=B_h(\Omega')$, and consequently $u_h$ tends to $u$ in the sense of $C^{\alpha}(\Omega')$ for every $\alpha'<\alpha$ and $\Omega' \Subset \Omega$, as $h\to 0$. Using Lemma 6.37 and Lemma 7.3 of the following section we can then conclude approximation results for $C^{k,\alpha}(\overline{\Omega})$ functions; (see Section 6.3).

We turn our attention now to the approximation of functions in the $L^p_{\mathrm{loc}}(\Omega)$ spaces.

\begin{lemma}\label{gt:lem-sobolev-regularization-lp}\dcref{Lemma 7.2}\group{ch07-sobolev-spaces}\source{gilbarg-trudinger:page-0160, gilbarg-trudinger:page-0161}
\textit{Let $u\in L^p_{\mathrm{loc}}(\Omega)(L^p(\Omega))$, $p<\infty$. Then $u_h$ converges to $u$ in the sense of $L^p_{\mathrm{loc}}(\Omega)(L^p(\Omega))$.}
\end{lemma}

\textit{Proof.} Using Hölder's inequality, we obtain from (7.13)

\[
\lvert u_h(x)\rvert^p\le \int_{|z|\le 1}\rho(z)\lvert u(x-hz)\rvert^p\,dz,
\]

% PDF page 161
\source{gilbarg-trudinger:page-0161}

so that if $\Omega' \Subset \Omega$ and $2h<\operatorname{dist}(\Omega',\partial\Omega)$,
\[
\int_{\Omega'} |u_h|^p\,dx \le \int_{\Omega'} \int_{|z|\le 1} \rho(z)|u(x-hz)|^p\,dz\,dx
\]
\[
= \int_{|z|\le 1} \rho(z)\,dz \int_{\Omega'} |u(x-hz)|^p\,dx
\]
\[
\le \int_{B_h(\Omega')} |u|^p\,dx,
\]
where $B_h(\Omega')=\{x \mid \operatorname{dist}(x,\Omega')<h\}$. Consequently
\begin{equation}
\tag{7.15}
\|u_h\|_{L^p(\Omega')} \le \|u\|_{L^p(\Omega'')}, \qquad \Omega''=B_h(\Omega').
\end{equation}
The proof can now be completed by approximation based on Lemma 7.1. Choose
$\varepsilon>0$ together with a $C^0(\Omega)$ function $w$ satisfying
\[
\|u-w\|_{L^p(\Omega'')} \le \varepsilon
\]
where $\Omega''=B_{h'}(\Omega')$ and $2h'<\operatorname{dist}(\Omega',\partial\Omega)$. By virtue of Lemma 7.1, we have for
sufficiently small $h$
\[
\|w-w_h\|_{L^p(\Omega'')} \le \varepsilon.
\]
Applying the estimate (7.15) to the difference $u-w$, we therefore obtain
\[
\|u-u_h\|_{L^p(\Omega')}
\le \|u-w\|_{L^p(\Omega')} + \|w-w_h\|_{L^p(\Omega')} + \|u_h-w_h\|_{L^p(\Omega')}
\]
\[
\le 2\varepsilon + \|u-w\|_{L^p(\Omega'')} \le 3\varepsilon
\]
for small enough $h\le h'$. Hence $u_h$ converges to $u$ in $L^p_{\mathrm{loc}}(\Omega)$. The result for $u\in L^p(\Omega)$
can then be obtained by extending $u$ to be zero outside $\Omega$ and applying the result
for $L^p_{\mathrm{loc}}(\mathbb{R}^n)$. $\square$

\section*{7.3. Weak Derivatives}

Let $u$ be locally integrable in $\Omega$ and $\alpha$ any multi-index. Then a locally integrable
function $v$ is called the $\alpha^{\text{th}}$ weak derivative of $u$ if it satisfies
\begin{equation}
\tag{7.16}
\int_\Omega \varphi v\,dx = (-1)^{|\alpha|} \int_\Omega u\,D^\alpha \varphi\,dx
\qquad \text{for all } \varphi \in C_0^{|\alpha|}(\Omega).
\end{equation}

% PDF page 162
\source{gilbarg-trudinger:page-0162}

We write $v=D^\alpha u$ and note that $D^\alpha u$ is uniquely determined up to sets of measure
zero. Pointwise relations involving weak derivatives will be accordingly understood
to hold almost everywhere. We call a function weakly differentiable if all its weak
derivatives of first order exist and $k$ times weakly differentiable if all its weak
derivatives exist for orders up to and including $k$. Let us denote the linear space of $k$
times weakly differentiable functions by $W^k(\Omega)$. Clearly $C^k(\Omega)\subset W^k(\Omega)$. The
concept of weak derivative is thus an extension of the classical concept which
maintains the validity of integration by parts (formula (7.16)).

We proceed to consider some basic properties of weakly differentiable functions.
The first lemma describes the interaction of weak derivatives and mollifiers.

\begin{lemma}\label{gt:lem-sobolev-mollifier-derivative}\dcref{Lemma 7.3}\group{ch07-sobolev-spaces}\source{gilbarg-trudinger:page-0162}
\emph{Let $u\in L^1_{\mathrm{loc}}(\Omega)$, $\alpha$ a multi-index, and suppose that $D^\alpha u$ exists. Then
if $\operatorname{dist}(x,\partial\Omega)>h$, we have}
\begin{equation}
\tag{7.17}
D^\alpha u_h(x)=(D^\alpha u)_h(x).
\end{equation}

\textbf{Proof.} \emph{By differentiating under the integral sign, we obtain}
\[
D^\alpha u_h(x)=h^{-n}\int_{\Omega} D_x^\alpha \rho\!\left(\frac{x-y}{h}\right)u(y)\,dy
\]
\[
=(-1)^{|\alpha|}h^{-n}\int_{\Omega} D^\alpha \rho\!\left(\frac{x-y}{h}\right)u(y)\,dy
\]
\[
=h^{-n}\int_{\Omega} \rho\!\left(\frac{x-y}{h}\right)D^\alpha u(y)\,dy \qquad \text{by (7.16)}
\]
\[
=(D^\alpha u)_h(x). \quad \square
\]

From Lemmas 7.1, 7.3 and the definition (7.16), now follows automatically a
basic approximation theorem for weak derivatives, the explicit verification of
which is left to the reader.

\end{lemma}

\begin{theorem}\label{gt:thm-sobolev-weak-derivative-approx}\dcref{Theorem 7.4}\group{ch07-sobolev-spaces}\source{gilbarg-trudinger:page-0162, gilbarg-trudinger:page-0163}
\emph{Let $u$ and $v$ be locally integrable in $\Omega$. Then $v=D^\alpha u$ if and only
if there exists a sequence of $C^\infty(\Omega)$ functions $\{u_m\}$ converging to $u$ in $L^1_{\mathrm{loc}}(\Omega)$ whose
derivatives $D^\alpha u_m$ converge to $v$ in $L^1_{\mathrm{loc}}(\Omega)$.}
\end{theorem}

This equivalent characterization of weak derivatives can also be used as their
definition, as is often the case. The resulting derivatives then are usually called
strong derivatives so that Theorem 7.4 establishes the equivalence of weak and
strong derivatives. Through Theorem 7.4, many results from the classical differential
calculus may be extended to weak derivatives simply by approximation. In
particular we have the \emph{product formula}

\begin{equation}
\tag{7.18}
D(uv)=uDv+vDu;
\end{equation}

% PDF page 163
\source{gilbarg-trudinger:page-0163}


\section*{7.4. The Chain Rule}

this holds for all $u, v \in W^1(\Omega)$ such that $uv,\, uDv + vDu \in \Loc(\Omega)$; (see Problem 7.4). Also if $\psi$ maps $\Omega$ onto a domain $\tilde{\Omega} \subset \R^n$ with $\psi \in C^1(\Omega)$, $\psi^{-1} \in C^1(\tilde{\Omega})$ and if $u \in W^1(\Omega)$, $v = u \circ \psi^{-1}$, then $v \in W^1(\tilde{\Omega})$ and the usual change of variables formula applies, that is

\[
(7.19)\qquad D_i u(x)=\frac{\partial y_j}{\partial x_i}\, D_{y_j} v(y)
\]

for almost all $x \in \Omega$, $y \in \tilde{\Omega}$, $y=\psi(x)$ (see Problem 7.5).
It is important to note that locally uniformly Lipschitz continuous functions are weakly differentiable, that is, $C^{0,1}(\Omega)\subset W^1(\Omega)$. This assertion follows since a function in $C^{0,1}(\Omega)$ will be absolutely continuous on any line segment in $\Omega$. Consequently its partial derivatives (which exist almost everywhere) satisfy (7.16) and hence coincide almost everywhere with the weak derivatives. By means of regularization, we can in fact prove that a function is weakly differentiable if and only if it is equivalent to a function that is absolutely continuous on almost all line segments in $\Omega$ parallel to the coordinate directions and whose partial derivatives are locally integrable; (see Problem 7.8). The basic properties of weak differentiation treated in this and the following section can be alternatively derived from this characterization.

\section*{7.4. The Chain Rule}

To complete our basic calculus of weak differentiation, we consider now a simple type of chain rule.

\begin{lemma}\label{gt:lem-sobolev-chain-rule-c1}\dcref{Lemma 7.5}\group{ch07-sobolev-spaces}\source{gilbarg-trudinger:page-0163, gilbarg-trudinger:page-0164}
\emph{Let $f \in C^1(\R)$, $f' \in L^\infty(\R)$ and $u \in W^1(\Omega)$. Then the composite function $f \circ u \in W^1(\Omega)$ and $D(f \circ u)=f'(u)\, Du$.}
\end{lemma}

\emph{Proof.} Let $u_m$, $m=1,2,\dots \in C^1(\Omega)$, and let $\{u_m\}$, $\{Du_m\}$ converge to $u$, $Du$ respectively in $L^1_{\Loc}(\Omega)$. Then for $\Omega' \Subset \Omega$, we have

\[
\int_{\Omega'} \lvert f(u_m)-f(u)\rvert\, dx \le \sup_{\Omega'} \lvert f'\rvert \int_{\Omega'} \lvert u_m-u\rvert\, dx \to 0 \qquad \text{as } m \to \infty
\]

\[
\int_{\Omega'} \lvert f'(u_m)Du_m-f'(u)Du\rvert\, dx \le \sup_{\Omega'} \lvert f'\rvert \int_{\Omega'} \lvert Du_m-Du\rvert\, dx
\]

\[
\qquad\qquad + \int_{\Omega'} \lvert f'(u_m)-f'(u)\rvert\, \lvert Du\rvert\, dx.
\]

A subsequence of $\{u_m\}$, which we renumber $\{u_m\}$, must converge a.e. $(\Omega')$ to $u$.
Since $f'$ is continuous, $\{f'(u_m)\}$ converges to $f'(u)$ a.e. $(\Omega')$. Hence the last integral

% PDF page 164
\source{gilbarg-trudinger:page-0164}

tends to zero by the dominated convergence theorem. Consequently the sequences
\[
\{f(u_m)\}, \{f'(u_m)Du_m\} \text{ tend to } f(u),\, f'(u)Du \text{ respectively, and therefore } Df(u)=
f'(u)Du.
\]
\begin{center}
$\square$
\end{center}

The positive and negative parts of a function $u$ are defined by
\[
u^+=\max\{u,0\}, \qquad u^-=\min\{u,0\}.
\]

Clearly $u=u^++u^-$ and $|u|=u^+-u^-$. From Lemma 7.5 we can derive the following
chain rule for these functions.

\begin{lemma}\label{gt:lem-sobolev-positive-negative-parts}\dcref{Lemma 7.6}\group{ch07-sobolev-spaces}\source{gilbarg-trudinger:page-0164}
\emph{Let $u\in W^{1}(\Omega)$; then $u^+$, $u^-$, $|u|\in W^{1}(\Omega)$ and}
\begin{equation}
\tag{7.20}
Du^+=
\begin{cases}
Du & \text{if } u>0 \\
0 & \text{if } u\le 0
\end{cases}
\qquad
Du^-=
\begin{cases}
0 & \text{if } u\ge 0 \\
Du & \text{if } u<0
\end{cases}
\qquad
D|u|=
\begin{cases}
Du & \text{if } u>0 \\
0 & \text{if } u=0 \\
-Du & \text{if } u<0.
\end{cases}
\end{equation}

\emph{Proof.} For $\varepsilon>0$, define
\[
f_\varepsilon(u)=
\begin{cases}
(u^2+\varepsilon^2)^{1/2}-\varepsilon & \text{if } u>0 \\
0 & \text{if } u\le 0.
\end{cases}
\]

Applying Lemma 7.5 we then have, for any $\varphi\in C_0^1(\Omega)$,
\[
\int_\Omega f_\varepsilon(u)D\varphi\,dx
=-\int_{u>0}\varphi\,\frac{uDu}{(u^2+\varepsilon^2)^{1/2}}\,dx
\]
and on letting $\varepsilon\to 0$, we obtain
\[
\int_\Omega u^+D\varphi\,dx
=-\int_{u>0}\varphi\,Du\,dx
\]
so that (7.20) is established for $u^+$. The other results follow since $u^-=-(-u)^+$
and $|u|=u^+-u^-$. $\square$

\end{lemma}

\begin{lemma}\label{gt:lem-sobolev-constant-gradient-zero}\dcref{Lemma 7.7}\group{ch07-sobolev-spaces}\source{gilbarg-trudinger:page-0164, gilbarg-trudinger:page-0165}
\emph{Let $u\in W^{1}(\Omega)$. Then $Du=0$ a.e. on any set where $u$ is constant.}
\end{lemma}

\emph{Proof.} Without loss of generality we may take the constant to be zero. The result
follows immediately from (7.20) since $Du=Du^++Du^-$. $\square$

% PDF page 165
\source{gilbarg-trudinger:page-0165}

We call a function piecewise smooth if it is continuous and has piecewise continuous first derivatives. The following chain rule then generalizes Lemmas 7.5 and 7.6.

\begin{theorem}\label{gt:thm-sobolev-piecewise-chain-rule}\dcref{Theorem 7.8}\group{ch07-sobolev-spaces}\source{gilbarg-trudinger:page-0165}
\emph{Let $f$ be a piecewise smooth function on $\mathbb{R}$ with $f' \in L^\infty(\mathbb{R})$. Then if $u \in W^{1}(\Omega)$, we have $f \circ u \in W^{1}(\Omega)$. Furthermore, letting $L$ denote the set of corner points of $f$, we have}
\begin{equation}
\tag{7.21}
D(f \circ u)=
\begin{cases}
f'(u)Du & \text{if } u \notin L \\
0 & \text{if } u \in L.
\end{cases}
\end{equation}

\emph{Proof.} By an induction argument the proof is reduced to the case of one corner which we may take without loss of generality at the origin. Let $f_1,f_2 \in C^{1}(\mathbb{R})$ satisfy $f_1',f_2' \in L^\infty(\mathbb{R})$, $f_1(u)=f(u)$ for $u \ge 0$, $f_2(u)=f(u)$ for $u \le 0$. Then since $f(u)=f_1(u^+)+f_2(u^-)$, the result follows by Lemmas 7.5 and 7.6. \hfill $\square$

Combining Lemma 7.7 and Theorem 7.8, we see that if $h$ is a finite valued function on $\mathbb{R}$, satisfying $h(u)=f'(u)$ for $u \notin L$, then $Df(u)=h(u)Du$. The chain rule in this form may be extended to Lipschitz continuous $f$ and $u \in W^{1}(\Omega)$ for which $h(u)Du \in L^1_{\mathrm{loc}}(\Omega)$. The proof of this assertion requires somewhat more measure theory than we have used; it is however a consequence of the characterization of weakly differentiable functions given in Problem 7.8.

\end{theorem}

\noindent\textbf{7.5. The $W^{k,p}$ Spaces}\label{gt:sec-wkp-spaces}\dcref{7.5}\group{ch07-sobolev-spaces}\source{gilbarg-trudinger:page-0165}

The $W^{k,p}(\Omega)$ spaces are Banach spaces analogous in a certain sense to the $C^{k,\alpha}(\overline{\Omega})$ spaces. In the $W^{k,p}(\Omega)$ spaces, continuous differentiability is replaced by weak differentiability and H\"older continuity by $p$-integrability. For $p \geq 1$ and $k$ a non-negative integer, we let
\[
W^{k,p}(\Omega)=\{u \in W^{k}(\Omega);\, D^\alpha u \in L^p(\Omega) \text{ for all } |\alpha| \leq k\}.
\]
The space $W^{k,p}(\Omega)$ is clearly linear. A norm is introduced by defining
\begin{equation}
\tag{7.22}
\|u\|_{k,p;\Omega}=\|u\|_{W^{k,p}(\Omega)}=
\left(\int_\Omega \sum_{|\alpha|\leq k} |D^\alpha u|^p\,dx\right)^{1/p}.
\end{equation}

We shall also use $\|u\|_{k,p}$ for $\|u\|_{k,p;\Omega}$ when there is no ambiguity. An equivalent norm would be
\begin{equation}
\tag{7.23}
\|u\|_{W^{k,p}(\Omega)}=\sum_{|\alpha|\leq k} \|D^\alpha u\|_p.
\end{equation}

The verification that $W^{k,p}(\Omega)$ is a Banach space under (7.22) is left to the reader (Problem 7.10).

% PDF page 166
\source{gilbarg-trudinger:page-0166}

Another Banach space $W_{0}^{k,p}(\Omega)$ arises by taking the closure of $C_{0}^{k}(\Omega)$ in
$W^{k,p}(\Omega)$. The spaces $W^{k,p}(\Omega)$, $W_{0}^{k,p}(\Omega)$ do not coincide for bounded $\Omega$. The case
$p=2$ is special, since the spaces $W^{k,2}(\Omega)$, $W_{0}^{k,2}(\Omega)$ (sometimes written $H^{k}(\Omega)$,
$H_{0}^{k}(\Omega)$) will be Hilbert spaces under the scalar product

\begin{equation}
\tag{7.24}
(u,v)_{k}=\int_{\Omega}\sum_{|\alpha|\le k} D^{\alpha}u\,D^{\alpha}v\,dx.
\end{equation}

Further functional analytic properties of $W^{k,p}(\Omega)$ and $W_{0}^{k,p}(\Omega)$ follow by con-
sidering their natural imbedding into the product of $N_{k}$ copies of $L^{p}(\Omega)$ where
$N_{k}$ is the number of multi-indices $\alpha$ satisfying $|\alpha|\le k$. Using the facts that finite
products and closed subspaces of separable (reflexive) Banach spaces are again
separable (reflexive) [DS], we obtain accordingly that the spaces $W^{k,p}(\Omega)$, $W_{0}^{k,p}(\Omega)$
are separable for $1\le p<\infty$ (reflexive for $1<p<\infty$).
The chain rule of Theorem 7.8 also extends to the spaces $W^{1,p}(\Omega)$, $W_{0}^{1,p}(\Omega)$.
In fact as a consequence of Theorem 7.8 and the definitions of these spaces we have
immediately that the space $W^{1}(\Omega)$ in the statement of Theorem 7.8 may be re-
placed by $W^{1,p}(\Omega)$, and by $W_{0}^{1,p}(\Omega)$ if also $f(0)=0$.
Local spaces $W^{k,p}_{\mathrm{loc}}(\Omega)$ can be defined to consist of functions belonging to
$W^{k,p}(\Omega')$ for all $\Omega'\Subset \Omega$. Theorem 7.4 shows that functions in $W^{k,p}_{\mathrm{loc}}(\Omega)$ with
compact support will in fact belong to $W_{0}^{k,p}(\Omega)$. Also, functions in $W^{1,p}(\Omega)$ which
vanish continuously on $\partial\Omega$ will belong to $W_{0}^{1,p}(\Omega)$, since they can be approximated
by functions with compact support.
In the case $p=\infty$, the Sobolev and Lipschitz spaces are related. In particular,
$W^{k,\infty}_{\mathrm{loc}}(\Omega)=C^{k-1,1}(\Omega)$ for arbitrary $\Omega$, and $W^{k,\infty}(\Omega)=C^{k-1,1}(\overline{\Omega})$ for sufficiently
smooth $\Omega$, e.g., for Lipschitz $\Omega$; (see Problem 7.7).

\subsection*{7.6. Density Theorems}

It is clear from Lemmas 7.2 and 7.3, that if $u$ lies in $W^{k,p}(\Omega)$, then $D^{\alpha}u_{h}$ tends to
$D^{\alpha}u$ in the sense of $L^{p}_{\mathrm{loc}}(\Omega)$ as $h$ approaches zero, for all multi-indices $\alpha$ satisfying
$|\alpha|\le k$. Using this fact we shall derive a global approximation result.

\begin{theorem}\label{gt:thm-sobolev-density-smooth}\dcref{Theorem 7.9}\group{ch07-sobolev-spaces}\source{gilbarg-trudinger:page-0166, gilbarg-trudinger:page-0167}
\emph{The subspace $C^{\infty}(\Omega)\cap W^{k,p}(\Omega)$ is dense in $W^{k,p}(\Omega)$.}
\end{theorem}

\emph{Proof.} Let $\Omega_{j}$, $j=1,2,\ldots$, be strictly contained subdomains of $\Omega$ satisfying
$\Omega_{j}\Subset \Omega_{j+1}$ and $\bigcup \Omega_{j}=\Omega$, and let $\{\psi_{j}\}$, $j=0,1,2,\ldots,$ be a partition of unity
(see Problem 6.8) subordinate to the covering $\{\Omega_{j+1}-\Omega_{j-1}\}$, $\Omega_{0}$ and $\Omega_{-1}$ being
defined as empty sets. Then for arbitrary $u\in W^{k,p}(\Omega)$ and $\varepsilon>0$, we can choose
$h_{j}$, $j=1,2,\ldots,$ satisfying

\begin{equation}
\tag{7.25}
h_{j}\le \operatorname{dist}(\Omega_{j},\partial\Omega_{j+1}), \qquad j\ge 1
\end{equation}

\[
\|(\psi_{j}u)_{h_{j}}-\psi_{j}u\|_{W^{k,p}(\Omega)}\le \frac{\varepsilon}{2^{j}}.
\]

% PDF page 167
\source{gilbarg-trudinger:page-0167}

Writing $v_j=(\psi_j u)_h$, we obtain from (7.25) that only a finite number of $v_j$ are non-vanishing on any given $\Omega' \Subset \Omega$. Consequently the function $v=\sum v_j$ belongs to $C^\infty(\Omega)$. Furthermore
\[
\|u-v\|_{W^{k,p}(\Omega)} \leq \sum \|v_j-\psi_j u\|_{W^{k,p}(\Omega)} \leq \varepsilon.
\]

This completes the proof. $\square$

Theorem 7.9 shows that $W^{k,p}(\Omega)$ could have been characterized as the completion of $C^\infty_0(\Omega)$ under the norm (7.22). In many instances this is a convenient definition.

In the case of arbitrary $\Omega$ we cannot replace $C^\infty(\Omega)$ by $C^\infty(\bar{\Omega})$ in Theorem 7.9. However, $C^\infty(\bar{\Omega})$ is dense in $W^{k,p}(\Omega)$ for a large class of domains $\Omega$ which includes for example $C^1$ domains (see Problem 7.11). More generally, if $\Omega$ satisfies a segment condition (that is, there exists a locally finite open covering $\{U_i\}$ of $\partial\Omega$ and corresponding vectors $y^i$ such that $x+ty^i \in \Omega$ for all $x \in \bar{\Omega}\cap U_i$, $t \in (0,1)$), then $C^\infty(\bar{\Omega})$ is dense in $W^{k,p}(\Omega)$. (See [AD].)

\subsection*{7.7. Imbedding Theorems}

This and the following section are concerned with the connection between pointwise and integrability properties of weakly differentiable functions and the integrability properties of their derivatives. One of the simplest results in this direction is that weakly differentiable functions of one variable must be absolutely continuous. In this section we prove the well known Sobolev inequalities for functions in $W^{1,p}_0(\Omega)$.

\begin{theorem}\label{gt:thm-sobolev-embedding-basic}\dcref{Theorem 7.10}\group{ch07-sobolev-spaces}\source{gilbarg-trudinger:page-0167, gilbarg-trudinger:page-0168, gilbarg-trudinger:page-0169, gilbarg-trudinger:page-0170}
\[
W^{1,p}_0(\Omega)=
\begin{cases}
L^{np/(n-p)}(\Omega) & \text{for } p<n \\
C^0(\bar{\Omega}) & \text{for } p>n.
\end{cases}
\]
Furthermore, there exists a constant $C=C(n,p)$ such that, for any $u\in W^{1,p}_0(\Omega)$,
\begin{equation}
\tag{7.26}
\|u\|_{np/(n-p)} \leq C\|Du\|_p \qquad \text{for } p<n,
\end{equation}
\[
\sup_{\Omega} |u| \leq C|\Omega|^{1/n-1/p}\|Du\|_p \qquad \text{for } p>n.
\]

\emph{Proof.} Let us first establish the estimates (7.26) for $C^1_0(\Omega)$ functions. We proceed from the case $p=1$. Clearly for any $u \in C^1_0(\Omega)$ and any $i$, $1 \leq i \leq n$,
\[
|u(x)| \leq \int_{-\infty}^{x_i} |D_i u|\, dx_i,
\]

% PDF page 168
\source{gilbarg-trudinger:page-0168}

\[
|u(x)|^{n/(n-1)} \le \left(\prod_{i=1}^n \int_{-\infty}^{\infty} |D_i u|\,dx_i\right)^{1/(n-1)} .
\]

The inequality (7.27) is now integrated successively over each variable $x_i$, $i=1,\ldots,n$, the generalized H\"older inequality (7.11) for $m=p_1=\cdots=p_m=n-1$ then being applied after each integration. Accordingly we obtain

\[
\|u\|_{n/(n-1)} \le \left(\prod_{i=1}^n \int_{\Omega} |D_i u|\,dx\right)^{1/n}
\]

\[
\le \frac{1}{n}\int_{\Omega}\sum_{i=1}^n |D_i u|\,dx
\]

\[
\le \frac{1}{\sqrt{n}}\|Du\|_1.
\]

Thus inequality (7.26) is established for the case $p=1$. The remaining cases can now be obtained by replacing $u$ in the estimate (7.28) by powers of $|u|$. In this way we get for $\gamma>1$,

\[
\||u|^\gamma\|_{n/(n-1)} \le \frac{\gamma}{\sqrt{n}} \int_{\Omega} |u|^{\gamma-1}|Du|\,dx
\]

\[
\le \frac{\gamma}{\sqrt{n}} \, \||u|^{\gamma-1}\|_{p'}\,\|Du\|_p
\]

by H\"older's inequality. Now for $p<n$ we may choose $\gamma$ to satisfy

\[
\frac{\gamma n}{n-1} = \frac{(\gamma-1)p}{p-1}, \qquad \text{i.e. } \gamma = \frac{(n-1)p}{n-p},
\]

and consequently obtain

\[
\|u\|_{np/(n-p)} \le \frac{\gamma}{\sqrt{n}}\|Du\|_p,
\]

as required.

The case $p>n$ follows immediately by combining inequalities (7.34), with $q=\infty$, $\mu=1/n$, and (7.37) of the following section. We insert here an alternative proof which is based on the case $p=1$.

For $p>n$, let us write

\[
\tilde{u} = \frac{\sqrt{n}\,|u|}{\|Du\|_p}
\]

% PDF page 169
\source{gilbarg-trudinger:page-0169}

and assume that $|\Omega| = 1$. We obtain then

$$
\| \tilde{u}^{\gamma} \|_{n'} \le \gamma \| \tilde{u}^{\gamma - 1} \|_{p'}, \qquad n' = \frac{n}{n-1}, \; p' = \frac{p}{p-1}
$$

so that

$$
\| \tilde{u} \|_{n'\gamma} \le \gamma^{1/\gamma} \| \tilde{u} \|_{p'(\gamma - 1)}^{1 - 1/\gamma}
\le \gamma^{1/\gamma} \| \tilde{u} \|_{p'}^{1 - 1/\gamma} \qquad \text{since } |\Omega| = 1 .
$$

Let us substitute for $\gamma$ the values $\delta^{\nu}$, $\nu = 1,2,\ldots$, where

$$
\delta = \frac{n'}{p'} > 1.
$$

We obtain thus

$$
\| \tilde{u} \|_{n'\delta^{\nu}} \le \delta^{\nu \delta^{-\nu}} \| \tilde{u} \|_{n'\delta^{\nu-1}}^{1 - \delta^{-\nu}}, \qquad \nu = 1,2,\ldots
$$

Iterating from $\nu = 1$ and using (7.28), we get for any $\nu$

$$
\| \tilde{u} \|_{\delta^{\nu}} \le \delta^{\sum \nu \delta^{-\nu}} \equiv \chi .
$$

Consequently as $\nu \to \infty$, we obtain by Problem 7.1,

$$
\sup_{\Omega} \tilde{u} \le \chi,
$$

and hence

$$
\sup_{\Omega} |u| \le \frac{\chi}{\sqrt{n}} \| Du \|_{p} .
$$

To eliminate the restriction $|\Omega| = 1$, we consider a transformation:
$y_i = |\Omega|^{1/n} x_i$.
We obtain thus

$$
\sup_{\Omega} |u| \le \frac{\chi}{\sqrt{n}} |\Omega|^{1/n - 1/p} \| Du \|_{p}
$$

as required.
To extend the estimates (7.26) to arbitrary $u \in W_0^{1,p}(\Omega)$, we let $\{u_m\}$ be a sequence of $C_0^1(\Omega)$ functions tending to $u$ in $W^{1,p}(\Omega)$. Applying the estimates (7.26)
to differences $u_{m_1} - u_{m_2}$, we see that $\{u_m\}$ will be a Cauchy sequence in $L^{np/(n-p)}(\Omega)$
for $p < n$ and in $C^0(\overline{\Omega})$ for $p > n$. Consequently the limit function $u$ will lie in the
desired spaces and satisfy (7.26). $\square$

% PDF page 170
\source{gilbarg-trudinger:page-0170}


Remark. The best constant $C$ satisfying (7.26) for the case $p<n$ was calculated
by Rodemich [RO], see also [BL], [TA 2], who showed that

$$
C=\frac{1}{n\sqrt{\pi}}
\left(
\frac{n!\,\Gamma(n/2)}{2\Gamma(n/p)\Gamma(n+1-n/p)}
\right)^{1/n}
\gamma^{1-1/p}, \qquad \gamma=\frac{n(p-1)}{n-p}.
$$

When $p=1$, the above number reduces to the well known isoperimetric constant
$n^{-1}(\omega_n)^{-1/n}$.

A Banach space $\B_1$ is said to be continuously imbedded in a Banach space
$\B_2$ (notation: $\B_1 \to \B_2$) if there exists a bounded, linear, one-to-one mapping:
$\B_1 \to \B_2$. Theorem 7.10 may be thus expressed as $W_0^{1,p}(\Omega)\to L^{np/(n-p)}(\Omega)$ if
$p<n$, $\to C^0(\bar{\Omega})$ if $p>n$. By iterating the result of Theorem 7.10 $k$ times we arrive at
an extension to the spaces $W_0^{k,p}(\Omega)$.

\medskip

\end{theorem}

\begin{corollary}\label{gt:cor-sobolev-embedding-higher-order}\dcref{Corollary 7.11}\group{ch07-sobolev-spaces}\source{gilbarg-trudinger:page-0170}
\textbf{Corollary 7.11.}

\[
W_0^{k,p}(\Omega)\to
\begin{cases}
L^{np/(n-kp)}(\Omega) & \text{for } kp<n,\\
C^m(\bar{\Omega}) & \text{for } 0\le m<k-\dfrac{n}{p}.
\end{cases}
\]

The second case is a consequence of the first, together with the case $p>n$ in
Theorem 7.10.

The estimates (7.26) and their extension to the spaces $W_0^{k,p}(\Omega)$ also show that a
norm on $W_0^{k,p}(\Omega)$ equivalent to (7.22) may be defined by

\begin{equation}
\tag{7.29}
\|u\|_{W_0^{k,p}(\Omega)}=
\left(
\int_{\Omega}\sum_{|\alpha|=k}|D^\alpha u|^p\,dx
\right)^{1/p}
\end{equation}

In general, $W_0^{k,p}(\Omega)$ cannot be replaced by $W^{k,p}(\Omega)$ in Corollary 7.11. However,
this replacement can be made for a large class of domains $\Omega$, which includes for
example domains with Lipschitz continuous boundaries. (See Theorem 7.26).
More generally, if $\Omega$ satisfies a uniform interior cone condition, (that is, there
exists a fixed cone $K_\Omega$ such that each $x\in\Omega$ is the vertex of a cone $K_\Omega(x)\subset\bar{\Omega}$
and congruent to $K_\Omega$), then there is an imbedding

\[
W^{k,p}(\Omega)\to
\begin{cases}
L^{np/(n-kp)}(\Omega) & \text{for } kp<n,\\
C_{B}^{m}(\Omega) & \text{for } 0\le m<k-\dfrac{n}{p}.
\end{cases}
\]

where $C_B^m(\Omega)=\{u\in C^m(\Omega)\mid D^\alpha u\in L^\infty(\Omega)\ \text{for }|\alpha|\le m\}$.

\end{corollary}

% PDF page 171
\source{gilbarg-trudinger:page-0171}

\begin{quote}
7.8. Potential Estimates and Imbedding Theorems
\end{quote}

\begin{center}
\noindent\textbf{7.8. Potential Estimates and Imbedding Theorems}\label{gt:sec-potential-estimates}\dcref{7.8}\group{ch07-sobolev-spaces}\source{gilbarg-trudinger:page-0171}
\end{center}

The imbedding results of the preceding section can be alternatively derived and
also improved through the use of certain potential estimates. Let $\mu \in (0,1]$
and define the operator $V_\mu$ on $L^1(\Omega)$ by the Riesz potential

\begin{equation}
(V_\mu f)(x)=\int_\Omega |x-y|^{n(\mu-1)}f(y)\,dy.
\tag{7.31}
\end{equation}

That $V_\mu$ is in fact well defined and maps $L^1(\Omega)$ into itself will appear as an incidental
consequence of the next lemma. First we observe, by setting $f \equiv 1$ in (7.31),

\begin{equation}
V_\mu 1 \le \mu^{-1}\omega_n^{1-\mu}|\Omega|^\mu.
\tag{7.32}
\end{equation}

For, choose $R>0$ so that $|\Omega|=|B_R(x)|=\omega_n R^n$. Then
\[
\int_\Omega |x-y|^{n(\mu-1)}\,dy \le \int_{B_R(x)} |x-y|^{n(\mu-1)}\,dy
\]
\[
=\mu^{-1}\omega_n R^{n\mu}
\]
\[
=\mu^{-1}\omega_n^{1-\mu}|\Omega|^\mu.
\]

\begin{lemma}\label{gt:lem-sobolev-riemann-potential-mapping}\dcref{Lemma 7.12}\group{ch07-sobolev-spaces}\source{gilbarg-trudinger:page-0171, gilbarg-trudinger:page-0172}
\emph{The operator $V_\mu$ maps $L^p(\Omega)$ continuously into $L^q(\Omega)$ for any $q$,
$1 \le q \le \infty$ satisfying}

\begin{equation}
0 \le \delta = \delta(p,q)=p^{-1}-q^{-1} < \mu.
\tag{7.33}
\end{equation}

\emph{Furthermore, for any $f \in L^p(\Omega)$,}

\begin{equation}
\|V_\mu f\|_q \le \left(\frac{1-\delta}{\mu-\delta}\right)^{1-\delta}\omega_n^{1-\mu}|\Omega|^{\mu-\delta}\|f\|_p.
\tag{7.34}
\end{equation}

\textit{Proof.} Choose $r \ge 1$ so that

\[
r^{-1}=1+q^{-1}-p^{-1}=1-\delta.
\]

Then it follows that $h(x-y)=|x-y|^{n(\mu-1)} \in L^r(\Omega)$, and by (7.32) one obtains

\[
\|h\|_r \le \left(\frac{1-\delta}{\mu-\delta}\right)^{1-\delta}\omega_n^{1-\mu}|\Omega|^{\mu-\delta}.
\]

The estimate (7.34) can now be derived by adapting the usual proof of the Young
inequality for convolutions in $\mathbb{R}^n$. Writing

% [illegible: final displayed factorization line]

% PDF page 172
\source{gilbarg-trudinger:page-0172}

we may estimate by the Hölder inequality (7.11)

\[
\lvert V_\mu f(x)\rvert \le
\left\{\int_\Omega h^r(x-y)\lvert f(y)\rvert^p\,dy\right\}^{1/q}
\left\{\int_\Omega h^r(x-y)\,dy\right\}^{1-1/p}
\left\{\int_\Omega \lvert f(y)\rvert^p\,dy\right\}^{\delta},
\]

so that

\[
\|V_\mu f\|_q \le \sup_\Omega \left\{\int_\Omega h^r(x-y)\,dy\right\}^{1/r}\,\|f\|_p
\le \left(\frac{1-\delta}{\mu-\delta}\right)^{1-\delta}\omega_n^{1-\mu}\lvert\Omega\rvert^{\mu-\delta}\,\|f\|_p.\quad \square
\]

We mention here that Lemma 7.12 may be strengthened in the sense that $V_\mu$
maps $L^p(\Omega)$ continuously into $L^q(\Omega)$ provided $p>1$ and
$\delta\le\mu$. The proof requires a well known integral inequality of Hardy and
Littlewood (see [HL]). However, Lemma 7.12 is adequate for our purposes here.
Observe that when $p>\mu^{-1}$, $V_\mu$ maps $L^p(\Omega)$ continuously into
$L^\infty(\Omega)$. Let us examine now the intermediate case
$p=\mu^{-1}$.

\medskip

\end{lemma}

\begin{lemma}\label{gt:lem-sobolev-riemann-potential-exponential}\dcref{Lemma 7.13}\group{ch07-sobolev-spaces}\source{gilbarg-trudinger:page-0172, gilbarg-trudinger:page-0173}
\textit{Let $f\in L^p(\Omega)$ and $g=V_{1/p}f$.
Then there exist constants $c_1$ and $c_2$ depending only on $n$ and $p$ such
that}

\[
\tag{7.35}
\int_\Omega \exp\!\left[\frac{g}{c_1\|f\|_p}\right]^{p'}\,dx \le c_2|\Omega|,
\qquad p' = p/(p-1).
\]

\noindent\textit{Proof.} From Lemma 7.12, we get for any $q\ge p$

\[
\|g\|_q \le q^{1-1/p+1/q}\omega_n^{1-1/p}|\Omega|^{1/q}\|f\|_p,
\]

\noindent so that

\[
\int_\Omega |g|^q\,dx \le q^{1+q/p'}\omega_n^{q/p'}|\Omega|\,\|f\|_p^q
\]

\noindent and hence for $q\ge p-1$

\[
\int_\Omega |g|^{p'q}\,dx \le p'q\bigl(\omega_n p'q\|f\|_p^{p'}\bigr)^q|\Omega|.
\]

% PDF page 173
\source{gilbarg-trudinger:page-0173}

Consequently

\[
\int_{\Omega}\sum_{N_0}^{N}\frac{1}{k!}\left(\frac{\abs{g}}{c_1\norm{f}_p}\right)^{p'k}\,dx
\le p'|\Omega|\sum\left(\frac{p'\omega_n}{c_1^{p'}}\right)^k\frac{k^k}{(k-1)!},
\qquad N_0=[p]
\]

The series on the right hand side converges provided $c_1^{p'}>e\omega_n p'$, whence by the
monotone convergence theorem and (7.8) the desired estimate (7.35) follows. $\square$

The next lemmas serve to clarify the connection between weak derivatives and
potentials of the above type.

\end{lemma}

\begin{lemma}\label{gt:lem-sobolev-newtonian-representation}\dcref{Lemma 7.14}\group{ch07-sobolev-spaces}\source{gilbarg-trudinger:page-0173, gilbarg-trudinger:page-0174}
\textit{Let $u\in W_0^{1,1}(\Omega)$. Then}

\begin{equation}\tag{7.36}
u(x)=\frac{1}{n\omega_n}\int_{\Omega}\frac{(x_i-y_i)D_i u(y)}{|x-y|^n}\,dy
\qquad \text{a.e. }(\Omega).
\end{equation}

\textit{Proof.}\quad Suppose that $u\in C_0^1(\Omega)$ and extend $u$ to be zero outside $\Omega$. Then, for any
$\omega$ with $|\omega|=1$,

\[
u(x)=-\int_0^{\infty}D_{\omega}u(x+r\omega)\,dr.
\]

Integrating with respect to $\omega$, we obtain

\[
u(x)=-\frac{1}{n\omega_n}\int_0^{\infty}\int_{|\omega|=1}D_{\omega}u(x+r\omega)\,dr\,d\omega
\]

\[
=\frac{1}{n\omega_n}\int_{\Omega}\frac{(x_i-y_i)D_i u(y)}{|x-y|^n}\,dy
\]

and (7.36) follows from Lemma 7.12 and the fact that $C_0^1(\Omega)$ is dense in
$W_0^{1,1}(\Omega)$. $\square$

Note that by means of the formula (7.16) for integration by parts, the Newtonian
potential representation for $C_0^2(\Omega)$ functions, equation (2.17), is deducible
from formula (7.36). Also we obtain for $u\in W_0^{1,1}(\Omega)$

\begin{equation}\tag{7.37}
|u|\le \frac{1}{n\omega_n}V_{1,n}|Du|.
\end{equation}

% PDF page 174
\source{gilbarg-trudinger:page-0174}

Combining Lemma 7.12 and inequality (7.37) we obtain immediately the embeddings
$W^{1,p}_0(\Omega)\to L^q(\Omega)$ for $p^{-1}-q^{-1}<n^{-1}$, which is almost the conclusion of Theorem 7.10. In fact, this weaker version would be adequate for the purposes of this book. But also combining Lemma 7.13 and (7.37), we obtain a sharpening of the case $p=n$ expressed by the following theorem.

\end{lemma}

\begin{theorem}\label{gt:thm-sobolev-moser-trudinger}\dcref{Theorem 7.15}\group{ch07-sobolev-spaces}\source{gilbarg-trudinger:page-0174}
\emph{Let $u\in W^{1,n}_0(\Omega)$. Then there exist constants $c_1$ and $c_2$ depending only on $n$, such that}

\begin{equation}
\tag{7.38}
\int_\Omega \exp\!\left(\frac{|u|}{c_1\|Du\|_n}\right)^{n/(n-1)}\,dx\leq c_2|\Omega|.
\end{equation}

\emph{Remark.} The estimate (7.37) is readily generalized to higher order weak derivatives. One obtains then for $u\in W^{k,1}_0(\Omega)$,

\begin{equation}
\tag{7.39}
|u|\leq \frac{1}{(k-1)! \, n \omega_n} V_{k,n}|D^k u|,
\end{equation}

and using Lemma 7.13 we have an extension of Theorem 7.15. Namely there exist constants $c_1$ and $c_2$ depending only on $n$ and $k$ such that if $u\in W^{k,p}_0(\Omega)$ with $n=kp$, then

\begin{equation}
\tag{7.40}
\int_\Omega \exp\!\left(\frac{|u|}{c_1\|D^k u\|_p}\right)^{p/(p-1)}\,dx\leq c_2|\Omega|.
\end{equation}

The case $p>n$ of the Sobolev imbedding theorem may be sharpened through the following lemma.

\end{theorem}

\begin{lemma}\label{gt:lem-sobolev-convex-average-estimate}\dcref{Lemma 7.16}\group{ch07-sobolev-spaces}\source{gilbarg-trudinger:page-0174, gilbarg-trudinger:page-0175}
\emph{Let $\Omega$ be convex and $u\in W^{1,1}(\Omega)$. Then}

\begin{equation}
\tag{7.41}
|u(x)-u_S|\leq \frac{d^n}{n|S|}\int_\Omega |x-y|^{1-n}|Du(y)|\,dy \qquad \text{a.e. }(\Omega),
\end{equation}

\emph{where}

\[
u_S=\frac{1}{|S|}\int_S u\,dx, \qquad d=\operatorname{diam}\Omega,
\]

\emph{and $S$ is any measurable subset of $\Omega$.}

\emph{Proof.} By Theorem 7.9, it is enough to establish (7.41) for $u\in C^1(\Omega)$. We then have for $x,y\in\Omega$,

\[
u(x)-u(y)=-\int_0^{|x-y|}D_{\omega}u(x+r\omega)\,dr,\qquad \omega=\frac{y-x}{|y-x|}.
\]

% PDF page 175
\source{gilbarg-trudinger:page-0175}

Integrating with respect to $y$ over $S$, we obtain

\[
|S|(u(x)-u_S)=-\int_S dy \int_0^{|x-y|} D_r u(x+r\omega)\,dr.
\]

Writing

\[
V(x)=
\begin{cases}
|D_r u(x)|, & x\in \Omega,\\
0, & x\notin \Omega
\end{cases}
\]

we thus have

\begin{align*}
|u(x)-u_S|
&\le \frac{1}{|S|}\int_{|x-y|<d} dy \int_0^\infty V(x+r\omega)\,dr \\
&= \frac{1}{|S|}\int_0^\infty \int_{|\omega|=1} \int_0^d V(x+r\omega)\rho^{n-1}\,d\rho\,d\omega\,dr \\
&= \frac{d^n}{n|S|}\int_0^\infty \int_{|\omega|=1} V(x+r\omega)\,d\omega\,dr \\
&= \frac{d^n}{n|S|}\int_\Omega |x-y|^{1-n}|D_r u(y)|\,dy. \quad \square
\end{align*}

We can now prove the embedding theorem of Morrey.

\end{lemma}

\begin{theorem}\label{gt:thm-sobolev-morrey}\dcref{Theorem 7.17}\group{ch07-sobolev-spaces}\source{gilbarg-trudinger:page-0175, gilbarg-trudinger:page-0176}
\emph{Let $u\in W^{1,p}_0(\Omega)$, $p>n$. Then $u\in C^\gamma(\overline{\Omega})$, where $\gamma=1-n/p$. Further-more, for any ball $B=B_R$,}

\[
(7.42)\qquad \osc_{\Omega\cap B} u \le CR^\gamma \|Du\|_p,
\]

\emph{where $C=C(n,p)$.}

\textit{Proof.} Coupling the estimates (7.41) and (7.34) for $S=\Omega=B$, $q=\infty$ and $\mu=n^{-1}$, we have

\[
|u(x)-u_B|\le C(n,p)R^\gamma \|Du\|_p \qquad \text{a.e. }(\Omega\cap B).
\]

The result then follows since

\[
|u(x)-u(y)|\le |u(x)-u_B|+|u(y)-u_B|
\le 2C(n,p)R^\gamma \|Du\|_p \qquad \text{a.e. }(\Omega\cap B). \quad \square
\]


% PDF page 176
\source{gilbarg-trudinger:page-0176}

Combining Theorems 7.10 and 7.17, we have for $u\in W_0^{1,p}(\Omega)$ and $p>n$ the estimate

\begin{equation}
\tag{7.43}
|u|_{0,\gamma}\le C[1+(\diam \Omega)^\gamma]\|Du\|_p.
\end{equation}

Further, the results of Theorems 7.10, 7.15, 7.17 may be summarized by the following diagram

\[
\begin{array}{c}
L^{np/(n-p)}(\Omega), \quad p<n \\[1.5ex]
W_0^{1,p}(\Omega)\to L^{\varphi}(\Omega), \qquad \varphi=\exp\!\bigl(|t|^{n/(n-1)}\bigr)-1,\quad p=n \\[1.5ex]
C^{\lambda}(\bar{\Omega}), \qquad \lambda=1-\frac{n}{p},\quad p>n
\end{array}
\]

where $L^{\varphi}(\Omega)$ denotes the Orlicz space with defining function $\varphi$. (See [TR 2] for a more explicit definition of $L^{\varphi}(\Omega)$.)

For the derivation of many of the a priori estimates in this book weaker forms of the Sobolev inequalities known as the Poincar\'e inequalities are sufficient. From Lemmas 7.12 and 7.14 we have for $u\in W_0^{1,p}(\Omega)$, $1\le p<\infty$

\begin{equation}
\tag{7.44}
\|u\|_p\le \left(\frac{1}{\omega_n}|\Omega|\right)^{1/n}\|Du\|_p;
\end{equation}

while from Lemmas 7.12 and 7.16 we have, for $u\in W^{1,p}(\Omega)$ and convex $\Omega$,

\begin{equation}
\tag{7.45}
\|u-u_S\|_p\le \left(\frac{\omega_n}{|S|}\right)^{1-1/n} d^n\|Du\|_p,\qquad d=\diam \Omega.
\end{equation}

\end{theorem}

\noindent\textbf{7.9. The Morrey and John-Nirenberg Estimates}\label{gt:sec-morrey-john-nirenberg}\dcref{7.9}\group{ch07-sobolev-spaces}\source{gilbarg-trudinger:page-0176}

We proceed now to a consideration of the potential operators $V_{\mu}$ on a different class of spaces in order to prove useful imbedding results due to Morrey (Theorem 7.19) and John and Nirenberg (Theorem 7.21). Namely, the integrable function $f$ is said to belong to $M^p(\Omega)$, $1\le p\le \infty$, if there exists a constant $K$ such that

\begin{equation}
\tag{7.46}
\int_{\Omega\cap B_R} |f|\,dx\le KR^{n(1-1/p)}
\end{equation}

for all balls $B_R$. We define the $p$ norm $\|f\|_{M^p(\Omega)}$ to be the infimum of the constants $K$ satisfying (7.46). It is easy to see that $L^p(\Omega)\subset M^p(\Omega)$, $L^1(\Omega)=M^1(\Omega)$, $L^{\infty}(\Omega)=M^{\infty}(\Omega)$. Instead of considering in detail the action of the operators $V_{\mu}$ on arbitrary $M^p(\Omega)$ spaces, it will be enough to limit ourselves to the cases $p\ge \mu^{-1}$.


% PDF page 177
\source{gilbarg-trudinger:page-0177}

\begin{quote}
\textbf{7.9. The Morrey and John-Nirenberg Estimates}\hfill 165
\end{quote}

\textbf{Lemma 7.18.} \emph{Let $f \in \M^p(\Omega)$, $\delta = p^{-1} < \mu$. Then}
\begin{equation}
|V_{\mu} f(x)| \le \frac{1-\delta}{\mu-\delta} (\diam \Omega)^{n(\mu-\delta)} \, \|f\|_{\M^p(\Omega)}
\qquad \text{a.e. }(\Omega).
\tag{7.47}
\end{equation}

\textbf{Proof.} Extend $f$ to be zero outside $\Omega$ and write
\[
v(\rho)=\int_{B_{\rho}(x)} |f(y)| \, dy.
\]
Then
\[
|V_{\mu} f(x)| \le \int_{\Omega} \rho^{n(\mu-1)} |f(y)| \, dy, \qquad \rho = |x-y|
\]
\[
= \int_{0}^{d} \rho^{n(\mu-1)} v'(\rho) \, d\rho, \qquad d = \diam \Omega
\]
\[
= d^{n(\mu-1)}v(d) + n(1-\mu)\int_{0}^{d} \rho^{n(\mu-1)-1} v(\rho)\, d\rho
\]
\[
\le \frac{1-\delta}{\mu-\delta} d^{n(\mu-\delta)} K \qquad \text{by (7.46).}\quad \square
\]

The following theorem now generalizes Theorem 7.17.

\textbf{Theorem 7.19.} \emph{Let $u \in W^{1,1}(\Omega)$, and suppose there exist positive constants $K$,
$\alpha$ $(\alpha \le 1)$ such that}
\begin{equation}
\int_{B_R} |Du| \, dx \le K R^{n-1+\alpha} \qquad \text{for all balls } B_R \subset \Omega.
\tag{7.48}
\end{equation}

\emph{Then $u \in C^{0,\alpha}(\Omega)$, and for any ball $B_R \subset \Omega$}
\begin{equation}
\osc_{B_R} u \le CKR^{\alpha},
\tag{7.49}
\end{equation}

\emph{where $C=C(n,\alpha)$. If $\Omega=\Omegae \cap \R^n_{+}=\{x\in \Omegae \mid x_n>0\}$ for some domain
$\Omegae \subset \R^n$ and (7.48) holds for all balls $B_R \subset \Omegae$, then $u \in C^{0,\alpha}(\Omegae\cap \Omega)$ and (7.49) holds for all
$B_R \subset \Omegae$.}

Theorem 7.19 is obtained by combining Lemma 7.16 ($S=\Omega$) with Lemma 7.18.
As a further consequence of Lemma 7.18 we have

\textbf{Lemma 7.20.} \emph{Let $f \in \M^p(\Omega)$ ($p>1$) and $g=V_{\mu}f$, $\mu=p^{-1}$. Then there exist constants $c_1$ and $c_2$ depending only on $n$ and $p$ such that}
\begin{equation}
\int_{\Omega} \exp\!\left(\frac{|g|}{c_1 K}\right) dx \le c_2 (\diam \Omega)^n
\tag{7.50}
\end{equation}

\emph{where $K=\|f\|_{\M^p(\Omega)}$.}

% PDF page 178
\source{gilbarg-trudinger:page-0178}

\emph{Proof.} Writing for any $q \geq 1$

\[
|x-y|^{n(\mu-1)} = |x-y|^{(\mu/q-1)n/q}\,|x-y|^{n(1-1/q)(\mu/q+\mu-1)}
\]

we have by H\"older's inequality

\[
|g(x)| \leq (V_{\mu/q}|f|)^{1/q}(V_{\mu+\mu/q}|f|)^{1-1/q}.
\]

\textbf{By Lemma 7.18}

\[
V_{\mu+\mu/q}|f| \leq \frac{(1-\mu)q}{\mu}\,d^{n/pq}K,\qquad d=\operatorname{diam}\Omega
\]

\[
\leq (p-1)q d^{n/pq}K.
\]

\textbf{Also by Lemma 7.12}

\[
\int_{\Omega} V_{\mu/q}|f|\,dx \leq pq\omega_n^{1-1/pq}|\Omega|^{1/pq}\|f\|_1
\]

\[
\leq pq\omega_n Kd^{n(1-1/p+1/pq)}.
\]

\textbf{Hence}

\[
\int_{\Omega} |g|^q\,dx \leq p(p-1)^{q-1}\omega_n^q d^n K^q
\]

\[
\leq p'\omega_n\{(p-1)qK\}^q d^n,\qquad p' = p/(p-1).
\]

\textbf{Consequently}

\[
\int_{\Omega}\sum_{m=0}^{N}\frac{|g|^m}{m!(c_1K)^m}\,dx
\leq p'\omega_n d^n \sum_{m=0}^{N}\left(\frac{p-1}{c_1}\right)^m \frac{m^m}{m!}
\]

\[
\leq c_2 d^n \qquad \text{if } (p-1)e < c_1.
\]

Letting $N \to \infty$, we thus obtain (7.50). $\square$

Combining Lemmas 7.16 and 7.20 we then get

\textbf{Theorem 7.21.} \emph{Let $u \in W^{1,1}(\Omega)$ where $\Omega$ is convex, and suppose there exists a constant $K$ such that}

\begin{equation}
\tag{7.51}
\int_{\Omega \cap B_R} |Du|\,dx \leq K R^{n-1} \qquad \text{for all balls } B_R.
\end{equation}

\emph{Then there exist positive constants $\sigma_0$ and $C$ depending only on $n$ such that}

\begin{equation}
\tag{7.52}
\int_{\Omega} \exp\!\left(\frac{\sigma}{K}|u-u_{\Omega}|\right)\,dx \leq C(\operatorname{diam}\Omega)^n
\end{equation}

\emph{where $\sigma = \sigma_0 |\Omega|(\operatorname{diam}\Omega)^{-n}$.}

% PDF page 179
\source{gilbarg-trudinger:page-0179}

\textbf{7.10. Compactness Results}

Let $\B_1$ be a Banach space continuously imbedded in a Banach space $\B_2$. Then
$\B_1$ is \emph{compactly imbedded} in $\B_2$ if the imbedding operator $I:\B_1 \to \B_2$ is compact,
that is, if the images of bounded sets in $\B_1$ are precompact in $\B_2$. Let us now prove
the \emph{Kondrachov compactness theorem} for the spaces $W^{1,p}_0(\Omega)$.

\medskip

\textbf{Theorem 7.22.} \emph{The spaces $W^{1,p}_0(\Omega)$ are compactly imbedded (i) in the spaces $L^q(\Omega)$
for any $q<np/(n-p)$, if $p<n$, and (ii) in $C^0(\overline{\Omega})$, if $p>n$.}

\medskip

\textit{Proof.} Part (ii) is a consequence of Morrey's theorem (Theorem 7.17) and
Arzela's theorem on equicontinuous families of functions. Let us thus concentrate
on part (i) and prove it initially for the case $q=1$. Let $A$ be a bounded set in $W^{1,p}_0(\Omega)$.
Without loss of generality we may assume that $A \subset C^1_0(\Omega)$ and that $\|u\|_{1,p;\Omega} \le 1$
for all $u \in A$. For $h>0$, we define $A_h=\{u_h \mid u \in A\}$ where $u_h$ is the regularization
of $u$ (see formula (7.13)). It then follows that the set $A_h$ is precompact in $L^1(\Omega)$.
For if $u \in A$ we have

\[
|u_h(x)| \le \int_{|z|\le 1} \rho(z)|u(x-hz)|\,dz \le h^{-n}\sup \rho \norm{u}_1
\]

and

\[
|Du_h(x)| \le h^{-1}\int_{|z|\le 1} |D\rho(z)||u(x-hz)|\,dz \le h^{-n-1}\sup |D\rho| \norm{u}_1
\]

so that $A_h$ is a bounded, equicontinuous subset of $C^0(\overline{\Omega})$ and hence precompact
in $C^0(\overline{\Omega})$ by Arzela's theorem, and consequently also precompact in $L^1(\Omega)$. Next
we may estimate for $u \in A$

\[
|u(x)-u_h(x)| \le \int_{|z|\le 1} \rho(z)|u(x)-u(x-hz)|\,dz
\]

\[
\le \int_{|z|\le 1} \rho(z)\int_0^{h|z|} |Du(x-r\omega)|\,dr\,dz,\qquad \omega=\frac{z}{|z|};
\]

hence integrating over $x$ we obtain

\[
\int_{\Omega}|u(x)-u_h(x)|\,dx \le h \int_{\Omega}|Du|\,dx \le h|\Omega|^{1-1/p}.
\]

Consequently $u_h$ is uniformly close to $u$ in $L^1(\Omega)$ (relative to $A$). Since we have
shown above that $A_h$ is totally bounded in $L^1(\Omega)$ for all $h>0$, it follows that $A$


% PDF page 180
\source{gilbarg-trudinger:page-0180}

is also totally bounded in $L^1(\Omega)$ and hence precompact. The case $q=1$ is thus established. To extend the result to arbitrary $q<np/(n-p)$, we estimate by (7.9)
\[
\|u\|_q\le \|u\|_1^{\lambda}\|u\|_{np/(n-p)}^{1-\lambda}
\qquad \text{where } \lambda+(1-\lambda)\left(\frac{1}{p}-\frac{1}{n}\right)=\frac{1}{q}
\]
\[
\le \|u\|_1^{\lambda}(C\|Du\|_p)^{1-\lambda}
\qquad \text{by Theorem 7.10.}
\]

Consequently a bounded set in $W_0^{1,p}(\Omega)$ must be precompact in $L^q(\Omega)$ for $q>1$
and the theorem is proved. $\square$

A simple extension of Theorem 7.22 shows that the imbeddings

\[
\begin{array}{c}
W_0^{k,p}(\Omega)\to L^q(\Omega)\qquad \text{for } kp<n,\ q<\dfrac{np}{n-kp}\\[2.2ex]
W_0^{k,p}(\Omega)\to C^m(\overline{\Omega})\qquad \text{for } 0\le m<k-\dfrac{n}{p}
\end{array}
\]

are compact and that $W_0^{k,p}(\Omega)$ may be replaced by $W^{k,p}(\Omega)$ for certain $\Omega$; see
Theorem 7.26; Problem 7.14.

\section*{7.11. Difference Quotients}

In partial differential equations, the weak or classical differentiability of functions
may often be deduced through a consideration of their difference quotients. Let $u$
be a function on a domain $\Omega$ in $\R^n$ and denote by $e_i$ the unit coordinate vector in
the $x_i$ direction. As in Chapter 6, we define the difference quotient in the direction
$e_i$ by

\begin{equation}
\tag{7.53}
\Delta^h u(x)=\Delta_i^h u(x)=\frac{u(x+he_i)-u(x)}{h},\qquad h\ne 0.
\end{equation}

The following basic lemmas pertain to difference quotients of functions in Sobolev
spaces.

\textbf{Lemma 7.23.} \emph{Let $u\in W^{1,p}(\Omega)$. Then $\Delta^h u\in L^p(\Omega')$ for any $\Omega'\subset \subset \Omega$ satisfying
$h<\dist(\Omega',\partial\Omega)$, and we have}

\[
\|\Delta^h u\|_{L^p(\Omega')}\le \|D_i u\|_{L^p(\Omega)}.
\]

\textbf{Proof.} \emph{Let us suppose initially that $u\in C^1(\Omega)\cap W^{1,p}(\Omega)$. Then}

\[
\Delta^h u(x)=\frac{u(x+he_i)-u(x)}{h}
\]

\[
=\frac{1}{h}\int_0^h \bigl[D_i u(x_1,\ldots,x_{i-1},x_i+\xi,x_{i+1},\ldots,x_n)\bigr]\,d\xi
\]

% PDF page 181
\source{gilbarg-trudinger:page-0181}

so that by H\"older's inequality

\[
|\Delta^h u(x)|^p \le \frac{1}{h}\int_0^h \left|D_i u(x_1,\ldots, x_{i-1}, x_i+\xi, x_{i+1},\ldots, x_n)\right|^p\,d\xi,
\]

and hence

\[
\int_{\Omega'} |\Delta^h u|^p\,dx \le \frac{1}{h}\int_0^h \int_{B_h(\Omega')} |D_i u|^p\,dx\,d\xi \le \int_{\Omega} |D_i u|^p\,dx.
\]

The extension to arbitrary functions in $W^{1,p}(\Omega)$ follows by a straight-forward
approximation argument using Theorem 7.9. $\square$

\medskip

\textbf{Lemma 7.24.} \emph{Let $u\in L^p(\Omega)$, $1<p<\infty$, and suppose there exists a constant $K$
such that $\Delta^h u\in L^p(\Omega')$ and $\|\Delta^h u\|_{L^p(\Omega')}\le K$ for all $h>0$ and $\Omega' \subset \subset \Omega$ satisfying
$h<\dist(\Omega',\partial\Omega)$. Then the weak derivative $D_i u$ exists and satisfies $\|D_i u\|_{L^p(\Omega)}\le K$.}

\medskip

\textit{Proof.} By the weak compactness of bounded sets in $L^p(\Omega')$, (Problem 5.4), there
exists a sequence $\{h_m\}$ tending to zero and a function $v\in L^p(\Omega)$ with $\|v\|_p\le K$
satisfying for all $\varphi\in C_0^1(\Omega)$

\[
\int_{\Omega} \varphi\, \Delta^{h_m}u\,dx \to \int_{\Omega} \varphi v\,dx.
\]

Now for $h_m<\dist(\operatorname{supp}\varphi,\partial\Omega)$, we have

\[
\int_{\Omega} \varphi\, \Delta^{h_m}u\,dx = -\int_{\Omega} u\, \Delta^{-h_m}\varphi\,dx \to -\int_{\Omega} u\, D_i\varphi\,dx.
\]

Hence

\[
\int_{\Omega} \varphi v\,dx = -\int_{\Omega} u\,D_i\varphi\,dx
\]

whence $v=D_i u$. $\square$

\bigskip

\textbf{7.12. Extension and Interpolation}

Under certain hypotheses on the domain $\Omega$, functions in Sobolev spaces $W^{k,p}(\Omega)$
may be extended as functions in $W^{k,p}(\R^n)$. We commence this section with a basic
extension result, analogous to Lemma 6.37, which will be used both to improve
previous imbedding results and to establish interpolation inequalities for Sobolev
space norms.


% PDF page 182
\source{gilbarg-trudinger:page-0182}

\textbf{Theorem 7.25.} \emph{Let $\Omega$ be a $C^{k-1,1}$ domain in $\R^n$, $k \geq 1$. Then (i) $C^\infty(\overline{\Omega})$ is dense in $W^{k,p}(\Omega)$, $1 \leq p < \infty$, and (ii) for any open set $\Omega' \supset \supset \Omega$ there exists a bounded linear extension operator $E$ from $W^{k,p}(\Omega)$ into $W_0^{k,p}(\Omega')$ such that $Eu = u$ in $\Omega$ and}

\begin{equation}
\tag{7.54}
\|Eu\|_{k,p;\Omega'} \leq C\|u\|_{k,p;\Omega}
\end{equation}

\emph{for all $u \in W^{k,p}(\Omega)$ where $C = C(k,\Omega,\Omega')$.}

\textbf{Proof.} We observe, by virtue of Lemmas 6.37 and 7.4, that assertions (i) and (ii) are equivalent. Let us first consider the density result (i) for the half-space $\R_+^n = \{x \in \R^n \mid x_n > 0\}$. In this case it is readily shown that the translated mollifications of $u$, given by

\begin{equation}
\tag{7.55}
v_h(x) = u_h(x + 2he_n)
\end{equation}

\[
= h^{-n}\int_{y_n > 0} u(y)\rho\!\left(\frac{x + 2he_n - y}{h}\right)\,dy,\qquad h > 0,
\]

converge to $u$ in $W^{k,p}(\R_+^n)$ as $h \to 0$. Accordingly an extension $E_0u$ of $u$ to all of $\R^n$ may be defined by the formula in Lemma 6.37, namely,

\begin{equation}
\tag{7.56}
E_0u(x) =
\begin{cases}
u(x) & \text{for } x_n > 0,\\
\sum_{i=1}^k c_i u(x',-x_n/i) & \text{for } x_n < 0
\end{cases}
\end{equation}

where $c_1,\ldots,c_k$ are constants determined by the system of equations

\[
\sum_{i=1}^k c_i(-1/i)^m = 1,\qquad m = 0,\ldots,k-1.
\]

If $u \in C^\infty(\R_+^n)\cap W^{k,p}(\R_+^n)$ it follows that $E_0u \in C^{k-1,1}(\R^n)\cap W^{k,p}(\R^n)$ and, moreover,

\begin{equation}
\tag{7.57}
\|E_0u\|_{k,p;\R^n} \leq C\|u\|_{k,p;\R_+^n},
\end{equation}

where $C = C(k)$. Therefore, by approximation we obtain that $E_0$ maps $W^{k,p}(\R_+^n)$ into $W^{k,p}(\R^n)$ and satisfies (7.57) for all $u \in W^{k,p}(\R_+^n)$.

Having treated the half-space case, let us now suppose that $\Omega$ is a $C^{k-1,1}$ domain in $\R^n$. According to the definition in Section 6.2, there exist a finite number of open sets $\Omega_j \subset \Omega'$, $j = 1,\ldots,N$, which cover $\partial\Omega$, and corresponding mappings $\psi_j$ of $\Omega_j$ onto the unit ball $B = B_1(0)$ in $\R^n$ such that

\[
\text{(i)}\quad \psi_j(\Omega_j\cap\Omega) = B^+ = B \cap \R_+^n;
\]
\[
\text{(ii)}\quad \psi_j(\Omega_j\cap\partial\Omega) = B \cap \partial\R_+^n;
\]
\[
\text{(iii)}\quad \psi_j \in C^{k-1,1}(\Omega_j),\qquad \psi_j^{-1}\in C^{k-1,1}(B).
\]

% PDF page 183
\source{gilbarg-trudinger:page-0183}

We let $\Omega_0 \subset \subset \Omega$ be a subdomain of $\Omega$ such that $\{\Omega_j\}, j = 0, \dots , N$, is a finite covering of $\Omega$, and let $\eta_j, j = 0, \dots , N$ be a partition of unity subordinate to this covering. Then $(\eta_j u) \circ \psi_j^{-1} \in W^{k,p}(\R^n_+)$ (Problem 7.5) and hence $E_0[(\eta_j u) \circ \psi_j^{-1}] \in W^{k,p}(\R^n)$, whence $E_0[(\eta_j u) \circ \psi_j^{-1}] \circ \psi_j \in W^{k,p}_0(\Omega_j), j = 1, \dots , N$, since $\operatorname{supp}\eta_j \subset \Omega_j$. Thus the mapping $E$ defined for $u \in W^{k,p}(\Omega)$ by

\begin{equation}
Eu = u\eta_0 + \sum_{j=1}^{N} E_0[(\eta_j u) \circ \psi_j^{-1}] \circ \psi_j
\tag{7.58}
\end{equation}

satisfies $Eu \in W^{k,p}_0(\Omega')$, $Eu = u$ in $\Omega$ and

\[
\|Eu\|_{k,p;\Omega'} \leq C \|u\|_{k,p;\Omega}
\]

where $C = C(k,N,\psi_j,\eta_j) = C(k,\Omega,\Omega')$. Furthermore $(Eu)_h \to u$ in $W^{k,p}(\Omega)$ as $h \to 0$. \hfill $\square$

By combining the case $k = 1$ in Theorem 7.25 with our previous imbedding results, Theorems 7.10, 7.12 and 7.22, we obtain corresponding imbedding results for the Sobolev spaces $W^{1,p}(\Omega)$ for Lipschitz domains $\Omega$. By iteration we then have the following general imbedding theorem for $W^{k,p}(\Omega)$.

\medskip

\noindent\textbf{Theorem 7.26.} \textit{Let $\Omega$ be a $C^{0,1}$ domain in $\R^n$. Then,}
\begin{enumerate}
\item[(i)] \textit{if $kp < n$, the space $W^{k,p}(\Omega)$ is continuously imbedded in $L^{p^*}(\Omega)$, $p^* = np/(n - kp)$, and compactly imbedded in $L^q(\Omega)$ for any $q < p^*$;}
\item[(ii)] \textit{if $0 \le m < k - \frac{n}{p} < m + 1$, the space $W^{k,p}(\Omega)$ is continuously imbedded in $C^{m,\alpha}(\Ol{\Omega})$, $\alpha = k - n/p - m$, and compactly imbedded in $C^{m,\beta}(\Ol{\Omega})$ for any $\beta < \alpha$.}
\end{enumerate}

We turn now to interpolation inequalities which we treat initially for the spaces $\Wokp(\Omega)$.

\medskip

\noindent\textbf{Theorem 7.27.} \textit{Let $u \in \Wokp(\Omega)$. Then for any $\epsilon > 0$, $0 < |\beta| < k$,}
\begin{equation}
\|D^\beta u\|_{p;\Omega} \le \epsilon \|u\|_{k,p;\Omega} + C \epsilon^{|\beta|/(|\beta|-k)} \|u\|_{p;\Omega},
\tag{7.59}
\end{equation}
\textit{where $C = C(k)$.}

\textit{Proof.} We establish (7.59) for the case $|\beta| = 1$, $k = 2$ which is needed in Chapter 9. A suitable induction argument yields the stated result for arbitrary $\beta, k$.

Let us first suppose $u \in C^2_0(\R)$ and consider an interval $(a,b)$ of length $b-a=\epsilon$. For $x' \in (a, a+\epsilon/3)$, $x'' \in (b-(\epsilon/3),b)$, we have, by the mean value theorem,

\[
|u'(\bar{x})| = \left| \frac{u(x') - u(x'')}{x' - x''} \right|
\]

\[
\le \frac{3}{\epsilon}\bigl(|u(x')| + |u(x'')|\bigr)
\]


% PDF page 184
\source{gilbarg-trudinger:page-0184}

for some $\bar{x}\in(a,b)$. Consequently for any $x\in(a,b)$,

\[
\lvert u'(x)\rvert \le \frac{3}{\varepsilon}\bigl(\lvert u(x')\rvert+\lvert u(x'')\rvert\bigr)+\int_a^b \lvert u''\rvert.
\]

Integrating with respect to $x'$ and $x''$, over the intervals $(a, a+\varepsilon/3)$, $(b-\varepsilon/3, b)$, respectively, we then obtain

\[
\lvert u'(x)\rvert \le \int_a^b \lvert u''\rvert+\frac{18}{\varepsilon^2}\int_a^b \lvert u\rvert,
\]

so that by H\"older's inequality

\[
\lvert u'(x)\rvert^p \le 2^{p-1}\left\{\varepsilon^{p-1}\int_a^b \lvert u''\rvert^p+\frac{(18)^p}{\varepsilon^{p+1}}\int_a^b \lvert u\rvert^p\right\}.
\]

Hence, integrating with respect to $x$ over $(a,b)$ we have

\[
\int_a^b \lvert u'(x)\rvert^p \le 2^{p-1}\left\{\varepsilon^p\int_a^b \lvert u''\rvert^p+\left(\frac{18}{\varepsilon}\right)^p\int_a^b \lvert u\rvert^p\right\}.
\]

Consequently if we subdivide $\mathbb{R}$ into intervals of length $\varepsilon$, we obtain by adding all such inequalities

\begin{equation}
\tag{7.60}
\int \lvert u'\rvert^p \le 2^{p-1}\left\{\varepsilon^p\int \lvert u''\rvert^p+\left(\frac{18}{\varepsilon}\right)^p\int \lvert u\rvert^p\right\}
\end{equation}

which is the desired result in the one-dimensional case. To extend to higher dimensions we fix $i$, $1\le i\le n$, and apply (7.60) to $u\in C_0^2(\Omega)$ regarded as a function of $x_i$ only. By successive integration over the remaining variables we thus obtain

\[
\int \lvert D_i u\rvert^p \le 2^{p-1}\left\{\varepsilon^p\int \lvert D_{ii}u\rvert^p+\left(\frac{18}{\varepsilon}\right)^p\int \lvert u\rvert^p\right\},
\]

so that

\[
\lVert D_i u\rVert_p \le \varepsilon \lVert D_{ii}u\rVert_p+\frac{C}{\varepsilon}\lVert u\rVert_p
\]

for $C=36$. \qed

\medskip

\textbf{By combining Theorems 7.25 and 7.27, we obtain interpolation inequalities for the Sobolev spaces $W^{k,p}(\Omega)$.}


% PDF page 185
\source{gilbarg-trudinger:page-0185}

\noindent\textbf{Theorem 7.28.} \quad Let $\Omega$ be a $C^{1,1}$ domain in $\R^n$ and $u \in W^{k,p}(\Omega)$. Then for any $\varepsilon > 0$, $0 < |\beta| < k$,

\begin{equation}
\tag{7.61}
\|D^\beta u\|_{p;\Omega} \le \varepsilon \|u\|_{k,p;\Omega} + C \varepsilon^{|\beta|/(|\beta|-k)}\|u\|_{p;\Omega}
\end{equation}

\noindent where $C = C(k,\Omega)$.

\medskip

Alternative derivations of interpolation inequalities are treated in Problems 2.15, 7.18 and 7.19. The density, extension, imbedding, and interpolation results of Theorems 7.25, 7.26 and 7.28 are all valid under less restrictive hypotheses on the domains $\Omega$; (see [AD]).

\bigskip

\textbf{Notes}

\medskip

For related material on Sobolev spaces the reader is referred to the books [AD], [FR], [MY 5] and [NE]. We have followed the custom of referring to the spaces of this chapter as Sobolev spaces although various notions of spaces of weakly differentiable functions were used prior to Sobolev's work [SO 1]; (in this regard see [MY 1] and [MY 5]). The process of mollification or regularization appeared in Friedrich's work [FD 1]. The density theorem, Theorem 7.9, is due to Meyers and Serrin [MS 2]. The Sobolev inequalities, Theorem 7.10, were essentially proved by Sobolev [SO 1, 2]; we have followed the proof of Nirenberg [NI 3] for the case $p < n$. The H\"older estimates, Theorems 7.17 and 7.19 were derived by Morrey [MY 1]. Theorem 7.21 is due to John and Nirenberg [JN]; our proof is taken from [TR 2] where also the estimate Theorem 7.15 appeared. The compactness result, Theorem 7.22, is due to Rellich [RE] in the case $p = 2$ and to Kondrachov [KN] for the general case.

\bigskip

\textbf{Problems}

\medskip

\noindent\textbf{7.1.} \quad Let $\Omega$ be a bounded domain in $\R^n$. If $u$ is a measurable function on $\Omega$ such that $|u|^p \in L^1(\Omega)$ for some $p \in \R$, we define

\[
\Phi_p(u)=\left[\frac{1}{|\Omega|}\int_\Omega |u|^p\,dx\right]^{1/p}.
\]

\noindent Show that: \quad (i) $\lim_{p\to+\infty}\Phi_p(u)=\sup_\Omega |u|$;

\hspace*{2.8em}(ii) $\lim_{p\to-\infty}\Phi_p(u)=\inf_\Omega |u|$;

\hspace*{2.8em}(iii) $\lim_{p\to 0}\Phi_p(u)=\exp\!\left[\frac{1}{|\Omega|}\int_\Omega \log |u|\,dx\right]$.

% PDF page 186
\source{gilbarg-trudinger:page-0186}



\setcounter{page}{174}
\noindent\textbf{7. Sobolev Spaces}\hfill 174

\bigskip

\noindent\textbf{7.2.} Show that a function $u$ is weakly differentiable in a domain $\Omega$ if and only if
it is weakly differentiable in a neighborhood of every point in $\Omega$.

\bigskip

\noindent\textbf{7.3.} Let $\alpha$, $\beta$ be multi-indices and $u$ be a locally integrable function on a domain
$\Omega$. Show that provided any two of the weak derivatives $D^{\alpha+\beta}u$, $D^\alpha(D^\beta u)$, $D^\beta(D^\alpha u)$
exist, they all exist and coincide a.e.\ $(\Omega)$.

\bigskip

\noindent\textbf{7.4.} Derive the product formula (7.18). (Hint: consider first the case, $u \in W^1(\Omega)$,
$v \in C^1(\Omega)$.)

\bigskip

\noindent\textbf{7.5.} Derive the formula (7.19), and show that it remains valid if we assume only
$\psi \in C^{0,1}(\Omega)$, $\psi^{-1} \in C^{0,1}(\bar{\Omega})$.

\bigskip

\noindent\textbf{7.6.} Let $\Omega$ be a domain in $\R^n$ containing the origin. Show that the function $\gamma$
given by $\gamma(x)=\lvert x\rvert^{-\alpha}$ belongs to $W^{k}(\Omega)$ provided $k+\alpha<n$.

\bigskip

\noindent\textbf{7.7.} Let $\Omega$ be a domain in $\R^n$. Show that a function $u \in C^{0,1}(\Omega)$ if and only if
$u$ is weakly differentiable with locally bounded weak derivatives.

\bigskip

\noindent\textbf{7.8.} Let $\Omega$ be a domain in $\R^n$. Show that a function $u$ is weakly differentiable in
$\Omega$ if and only if it is equivalent to a function $\bar{u}$ that is absolutely continuous on
almost all line segments in $\Omega$ parallel to the coordinate axes and whose partial
derivatives, (which consequently exist a.e.\ $(\Omega)$), are locally integrable in $\Omega$. (See
[MY 5], p.\ 66). Derive from this characterization the product formula and chain
rule for weak differentiation.

\bigskip

\noindent\textbf{7.9.} Show that the norms (7.22) and (7.23) are equivalent norms on $W^{k,p}(\Omega)$.

\bigskip

\noindent\textbf{7.10.} Prove that the space $W^{k,p}(\Omega)$ is complete under either of the norms (7.22),
(7.23).

\bigskip

\noindent\textbf{7.11.} Let $\Omega$ be a domain whose boundary can be locally represented as the graph
of a Lipschitz continuous function. Show that $C^\infty(\bar{\Omega})$ is dense in $W^{k,p}(\Omega)$ for
$1 \leq p < \infty$, $k \geq 1$, and compare this result with the density result in Theorem 7.25.

\bigskip

\noindent\textbf{7.12.} Let $\Omega$ be a $C^{0,1}$ domain. For any function $u \in W^{1,p}(\Omega)$ and $1 \leq p < n$, derive
the Sobolev--Poincar\'e inequality

\[
\norm{u-u_\Omega}{\frac{np}{n-p};\Omega} \leq C\norm{Du}{p;\Omega}
\]

\noindent(where $C$ is independent of $u$) by a contradiction argument based on the compactness
result of Theorem 7.26.

\bigskip

\noindent\textbf{7.13.} Deduce from Theorem 7.19 the corresponding global result. Namely let
$u \in W^1(\Omega)$, $\partial \Omega \in C^{0,1}$ and suppose there exist positive constants $K$, $\alpha$ $(\alpha < 1)$ such
that

\[
\int_{B_R} \lvert Du\rvert\, dx \leq KR^{n-1+\alpha}\qquad \text{for all balls } B_R \subset \R^n.
\]


% PDF page 187
\source{gilbarg-trudinger:page-0187}



\setcounter{page}{175}
\noindent\textbf{Problems}\hfill 175

\bigskip

\noindent Then $u \in C^{0,\alpha}(\bar{\Omega})$ and

\[
[u]_{\alpha;\Omega} \leq CK
\]

\noindent where $C = C(n,\alpha,\Omega)$.

\bigskip

\noindent\textbf{7.14.} \ Let $\Omega$ be a bounded domain for which an imbedding

\[
W^{1,p}(\Omega)\to L^{p^*}(\Omega), \qquad 1 \leq p < \infty,
\]

\noindent is valid. Show that the imbedding

\[
W^{1,p}(\Omega)\to L^q(\Omega)
\]

\noindent is compact for any $q < p^*$.

\bigskip

\noindent\textbf{7.15.} \ Let $\Omega$ be a domain in $\R^n$. The total variation of a function $u \in L^1(\Omega)$ is
defined by

\[
\int_\Omega \abs{Du} = \sup \left\{ \int_\Omega u \,\operatorname{div} v \mid v \in \C_0^1(\Omega), \abs{v} \leq 1 \right\}.
\]

\noindent Show that the space $\BV(\Omega)$ of functions of finite total variation is a Banach space
under the norm

\[
\norm{u}{\BV(\Omega)} = \norm{u}{1} + \int_\Omega \abs{Du},
\]

\noindent and that $W^{1,1}(\Omega)$ is a closed subspace.

\bigskip

\noindent\textbf{7.16.} \ Let $u \in \BV(\Omega)$. By invoking the regularization of $u$ and appropriately
modifying the proof of Theorem 7.9, show that there exists a sequence $\{u_m\} \subset$
$\C^\infty(\Omega) \cap W^{1,1}(\Omega)$ such that $u_m \to u$ in $L^1(\Omega)$ and

\[
\int_\Omega \abs{Du_m} \to \int_\Omega \abs{Du}.
\]

\bigskip

\noindent\textbf{7.17} \ Let $\Omega$ be a bounded domain for which the Sobolev imbedding

\[
W^{1,1}(\Omega)\to L^{n/(n-1)}(\Omega)
\]

\noindent is valid. Show that also

\[
\BV(\Omega)\to L^{n/(n-1)}(\Omega)
\]


% PDF page 188
\source{gilbarg-trudinger:page-0188}

and furthermore that the embedding

\[
BV(\Omega)\to L^q(\Omega)
\]

is compact for any $q<n/(n-1)$.

\begin{enumerate}
\item[7.18.] Derive Theorem 7.27 for $p \geq 2$ from Green's first identity (2.10); (see Problem 2.15).

\item[7.19.] Let $\Omega$ be a $C^{0,1}$ domain. Derive the interpolation inequality (7.61) in the weaker form

\[
\|D^\beta u\|_{p;\Omega} \leq \epsilon \|u\|_{k,p;\Omega} + C_\epsilon \|u\|_{p;\Omega},
\]

\noindent (C$_\epsilon$ independent of $u$), by means of a contradiction argument based on the compactness result of Theorem 7.26.

\item[7.20.] Using regularization, show that locally integrable solutions of Laplace's equation (in the sense of Problem 2.8) are smooth and hence deduce the validity of the interior estimates in Chapter 4 for such solutions of Poisson's equation.

\item[7.21.] Using Morrey's inequality (7.42), prove that functions in the Sobolev space $W^{1,p}(\Omega)$, where $p>n$, are classically differentiable almost everywhere in $\Omega$.
\end{enumerate}


% PDF page 189
\source{gilbarg-trudinger:page-0189}
\pagestyle{empty}

Chapter 8

Generalized Solutions and Regularity

This chapter treats linear elliptic operators having principal part in divergence
form under relatively weak smoothness assumptions on the coefficients. We
consider operators $L$ of the form

\begin{equation}
Lu = D_i(a^{ij}(x)D_j u + b^i(x)u) + c^i(x)D_i u + d(x)u
\tag{8.1}
\end{equation}

whose coefficients $a^{ij}$, $b^i$, $c^i$, $d$ $(i,j = 1, . . . , n)$ are assumed to be measurable
functions on a domain $\Omega \subset \mathbb{R}^n$. An operator $L$ of the general form (3.1) may be
written in the form (8.1) provided its principal coefficients $a^{ij}$ are differentiable.
The Hilbert space approach developed here can then be viewed as providing an
alternative existence theory to that of Chapter 6. On the other hand, if in (8.1) the
coefficients $a^{ij}$ and $b^i$ are differentiable and the function $u \in C^2(\Omega)$, then $L$ may be
written in the general form (3.1) so that the theory of Chapter 6 would apply. The
divergence form however has the advantage that the operator $L$ may be defined for
significantly broader classes of functions than the class $C^2(\Omega)$. Indeed, if we assume
that the function $u$ is only weakly differentiable and that the functions $a^{ij}D_j u + b^i u$,
$c^i D_i u + du$, $i = 1, . . . , n$ are locally integrable, then, in a weak or generalized sense,
$u$ is said to satisfy $Lu = 0$ ($\geq 0$, $\leq 0$) respectively in $\Omega$ according as

\begin{equation}
\mathcal{L}(u, v) = \int_\Omega \{(a^{ij}D_j u + b^i u)D_i v - (c^i D_i u + du)v\}\,dx = 0 \, (\leq 0, \geq 0)
\tag{8.2}
\end{equation}

for all non-negative functions $v \in C_0^1(\Omega)$. Provided the coefficients of $L$ are locally
integrable, it follows from the divergence theorem (2.3) that a function $u \in C^2(\Omega)$
satisfying $Lu = 0$ ($\geq 0$, $\leq 0$) in the classical sense also satisfies these relations in the
generalized sense. Moreover, if the coefficients $a^{ij}$, $b^i$ have locally integrable
derivatives, then a generalized solution $u \in C^2(\Omega)$ is also a classical solution.

Let $f^i$, $g$, $i = 1, . . . , n$ be locally integrable functions in $\Omega$. Then a weakly
differentiable function $u$ will be called a weak or generalized solution of the
inhomogeneous equation

\begin{equation}
Lu = g + D_i f^i
\tag{8.3}
\end{equation}

in $\Omega$ if
